\section{Sistemas Operacionais}

Um sistema operacional é um software responsável por gerenciar os recursos de hardware de um sistema computacional e fornecer serviços para a execução de aplicações.
Atuando como uma camada de abstração entre o hardware e os programas, o sistema operacional simplifica o desenvolvimento de software ao oferecer mecanismos padronizados
para acessar o processador, memória, dispositivos de entrada e saída, dispositivos armazenamento de dados, etc.

A evolução dos sistemas operacionais permitiu o surgimento de plataformas cada vez mais complexas, capazes de executar múltiplas aplicações simultaneamente e
fornecer diversos serviços ao usuário. Entretanto, dependendo dos requisitos da aplicação, características como previsibilidade temporal, confiabilidade e isolamento podem
se tornar tão importantes quanto o desempenho, motivando o desenvolvimento de sistemas operacionais especializados.
